\section*{Capítulo \ref{ch:g8:equations} --- Cálculo literal y ecuaciones}

\begin{solution}{exo:g8:equations:1}
$5(2x + 3) = 10x + 15$.

$-3(4x - 1) = -12x + 3$.

$2(3x - 2) + 4(x + 5) = 6x - 4 + 4x + 20 = 10x + 16$.
\end{solution}

\begin{solution}{exo:g8:equations:2}
$(x+2)(x+7) = x^2 + 7x + 2x + 14 = x^2 + 9x + 14$.

$(3x+1)(2x-5) = 6x^2 - 15x + 2x - 5 = 6x^2 - 13x - 5$.

$(x-4)(x-6) = x^2 - 6x - 4x + 24 = x^2 - 10x + 24$.
\end{solution}

\begin{solution}{exo:g8:equations:3}
$4 \times 3 - 5 = 7$: sí. \quad
$2 \times 3 + 1 = 7$ y $3 + 4 = 7$: sí. \quad
$3^2 = 9$ y $6 \times 3 - 9 = 9$: sí --- $x = 3$ resuelve las
tres.
\end{solution}

\begin{solution}{exo:g8:equations:4}
$x + 9 = 4$: $x = -5$ (comprobación: $-5 + 9 = 4$).

$5x = -35$: $x = -7$.

$\frac x3 = 8$: $x = 24$.

$2x - 7 = 13$: $2x = 20$, $x = 10$ (comprobación: $20 - 7 = 13$).
\end{solution}

\begin{solution}{exo:g8:equations:5}
$6x + 5 = 2x + 25$: $4x = 20$, luego $x = 5$.

$9 - 4x = x - 11$: $20 = 5x$, luego $x = 4$.

$3(x + 2) = 2x + 11$: $3x + 6 = 2x + 11$, luego $x = 5$.
\end{solution}

\begin{solution}{exo:g8:equations:6}
$8x - 12 = 5x + 30$, luego $3x = 42$ y $x = 14$. Comprobación:
$4(28 - 3) = 100$ y $5(14 + 6) = 100$.
\end{solution}

\begin{solution}{exo:g8:equations:7}
``$3x - 8 = 19$'': $3x = 27$, $x = 9$.

``$2x + 5 = x + 12$'': $x = 7$.
\end{solution}

\begin{solution}{exo:g8:equations:8}
Ancho $x$, largo $x + 5$; \omterm{def:g6:measure:perimeter}{perímetro}
$2\bigl(x + (x + 5)\bigr) = 4x + 10 = 58$, luego $4x = 48$ y $x = 12$:
ancho $12$ cm, largo $17$ cm (comprobación:
$2 \times (12 + 17) = 58$).
\end{solution}

\begin{solution}{exo:g8:equations:9}
Los números son $n - 1$, $n$, $n + 1$, de \omterm{def:g1:addition:def}{suma} $3n = 141$, luego
$n = 47$: los números son $46$, $47$, $48$.
\end{solution}

\begin{solution}{exo:g8:equations:10}
$3x + 4 > 19$: $3x > 15$, luego $x > 5$.

$7 - 2x \geq 1$: $-2x \geq -6$; al dividir entre $-2$ se invierte:
$x \leq 3$.

$5(x-1) < 2x + 7$: $5x - 5 < 2x + 7$, luego $3x < 12$ y $x < 4$.
\end{solution}

\begin{solution}{exo:g8:equations:11}
Al cabo de $n$ meses, Mía tiene $140 + 15n$ y Tomás $200 + 10n$. Mía
va por delante cuando
\[
140 + 15n > 200 + 10n
\ \Longleftrightarrow\
5n > 60
\ \Longleftrightarrow\
n > 12 :
\]
a partir del mes $13$.
\end{solution}

\begin{solution}{exo:g8:equations:12}
Multiplica los dos miembros por $20$:
\[
5(x + 3) = 4(2x - 1)
\ \Longrightarrow\
5x + 15 = 8x - 4
\ \Longrightarrow\
19 = 3x
\ \Longrightarrow\
x = \frac{19}{3}.
\]
Comprobación: $\frac{19/3 + 3}{4} = \frac{28/3}{4} = \frac{28}{12} =
\frac73$, y $\frac{2 \times 19/3 - 1}{5} = \frac{35/3}{5} =
\frac{35}{15} = \frac73$. Iguales: la solución es $\frac{19}{3}$.
\end{solution}

\begin{solution}{pb:g8:equations:1}
\textbf{1.} Doble distributividad
(\cref{thm:g8:equations:expand}) con $c = a$, $d = b$:
\[
(a + b)(a + b) = a^2 + ab + ba + b^2 = a^2 + 2ab + b^2 ,
\]
porque $ab$ y $ba$ son el mismo \omterm{def:g2:mult:def}{producto}.

\textbf{2.} La misma expansión, con cuidado con los signos:
\[
(a - b)(a - b) = a^2 - ab - ba + b^2 = a^2 - 2ab + b^2 ,
\]
\[
(a + b)(a - b) = a^2 - ab + ba - b^2 = a^2 - b^2 :
\]
esta vez los términos cruzados $-ab$ y $+ba$ se cancelan en lugar de
duplicarse.

\textbf{3.} El cuadrado de lado $a + b$ se recorta en cuatro piezas: un
cuadrado $a \times a$ (el $a^2$), un cuadrado $b \times b$ (el
$b^2$) y \emph{dos} rectángulos $a \times b$ (juntos, el
$2ab$). La identidad $(a+b)^2 = a^2 + 2ab + b^2$ dice justamente
que las cuatro \omterm{def:g6:measure:area}{áreas} llenan el cuadrado grande --- es el dibujo de
la doble distributividad con los dos lados partidos igual.

\textbf{4.} Con las identidades:
\begin{align*}
101^2 &= (100 + 1)^2 = 10\,000 + 2 \times 100 + 1 = 10\,201, \\
99^2 &= (100 - 1)^2 = 10\,000 - 200 + 1 = 9\,801, \\
49 \times 51 &= (50 - 1)(50 + 1) = 50^2 - 1^2 = 2\,499 .
\end{align*}

\textbf{5.} Para $a = 3$, $b = 4$: $(3 + 4)^2 = 49$, pero
$3^2 + 4^2 = 9 + 16 = 25$. La fórmula de Zoe olvida los dos
rectángulos $a \times b$ del dibujo --- faltan los
$2ab = 24$, y en efecto $25 + 24 = 49$.

\textbf{6.} $1$; $1 + 3 = 4$; $1 + 3 + 5 = 9$;
$1 + 3 + 5 + 7 = 16$; $1 + 3 + 5 + 7 + 9 = 25$. Los resultados son
los cuadrados perfectos $1^2, 2^2, 3^2, 4^2, 5^2$.

\textbf{7.} Los \omterm{def:g3:numbers:evenodd}{números impares} empiezan en $1$ y avanzan de $2$ en
$2$: al $n$-ésimo se llega tras $n - 1$ saltos de $2$ desde
$1$, así que vale $1 + 2(n - 1) = 2n - 1$. Comprobación: $n = 1, 2, 3$
dan $1, 3, 5$.

\textbf{8.} Por la primera identidad,
\[
(n + 1)^2 - n^2 = n^2 + 2n + 1 - n^2 = 2n + 1 :
\]
la \omterm{ex:g1:subtraction:difference}{diferencia} entre dos cuadrados consecutivos es exactamente el
\omterm{def:g3:numbers:evenodd}{número impar} $(n+1)$-ésimo (pregunta 7 con $n + 1$ en lugar de $n$:
$2(n+1) - 1 = 2n + 1$).

\textbf{9.} Construye los cuadrados uno tras otro. Empieza con una
sola casilla: $1 = 1^2$. Para pasar del cuadrado de lado $1$ al
cuadrado de lado $2$, añade $2 \times 1 + 1 = 3$ casillas (pregunta
8): $1 + 3 = 2^2$. Para pasar al lado $3$, añade $2 \times 2 + 1 = 5$
casillas: $1 + 3 + 5 = 3^2$. Cada nuevo \omterm{def:g3:numbers:evenodd}{número impar} es exactamente la
capa en forma de L que completa el cuadrado siguiente, así que al añadir los
$n$ primeros \omterm{def:g3:numbers:evenodd}{números impares} el cuadrado de lado $n$ queda completo:
$1 + 3 + \dots + (2n - 1) = n^2$.

\textbf{10.} Los \omterm{def:g3:numbers:evenodd}{números impares} de $1$ a $99$ son los $50$ primeros
($2n - 1 = 99$ da $n = 50$), luego
\[
1 + 3 + 5 + \dots + 99 = 50^2 = 2\,500 .
\]
Y $1 + 3 + \dots + (2n - 1) = 400 = 20^2$ exige $n = 20$:
los \emph{veinte} primeros \omterm{def:g3:numbers:evenodd}{números impares} ($1$ hasta $39$) suman
$400$.

\textbf{11.} $x^2 + 6x + 9 = x^2 + 2 \times 3 \times x + 3^2 =
(x + 3)^2$: la primera identidad con $a = x$, $b = 3$. Así que la
\omterm{def:g8:equations:def}{ecuación} se lee $(x + 3)^2 = 0$. Un cuadrado es \omterm{def:g1:counting:zero}{cero} solo cuando el
número mismo es \omterm{def:g1:counting:zero}{cero} (\cref{thm:g8:negprod:rules}: un número no
nulo tiene distancia no nula, luego su cuadrado también tiene distancia
no nula). Por tanto $x + 3 = 0$: la única solución es $x = -3$.

\textbf{12.} Por el \cref{thm:g8:negprod:rules}, la distancia de un
\omterm{def:g2:mult:def}{producto} es el \omterm{def:g2:mult:def}{producto} de las distancias. Si los dos factores son
no nulos, las dos distancias son no nulas, luego la distancia del \omterm{def:g2:mult:def}{producto} es
no nula: el \omterm{def:g2:mult:def}{producto} no puede ser $0$. Un \omterm{def:g2:mult:def}{producto} \omterm{def:g1:counting:zero}{cero} obliga, pues,
a que al menos un factor sea \omterm{def:g1:counting:zero}{cero}. Para $(x - 2)(x + 5) = 0$: o bien
$x - 2 = 0$, o bien $x + 5 = 0$, lo que da las dos soluciones $x = 2$ y
$x = -5$.

\textbf{13.} $x^2 = 49$ se convierte en $x^2 - 49 = 0$, y la tercera
identidad lo factoriza:
\[
x^2 - 49 = x^2 - 7^2 = (x - 7)(x + 7) = 0 .
\]
Por la pregunta 12, $x = 7$ o $x = -7$: \emph{dos} soluciones. Todas las
ecuaciones resueltas hasta aquí en este capítulo eran de primer grado y
tenían exactamente una solución; un cuadrado en la \omterm{def:g8:equations:def}{ecuación} puede dar
dos.

\textbf{14.} Pasa todo al miembro de la izquierda:
$x^2 - 6x + 9 = 0$, y reconoce la segunda identidad:
$x^2 - 2 \times 3 \times x + 3^2 = (x - 3)^2$. Así que la \omterm{def:g8:equations:def}{ecuación} es
$(x - 3)^2 = 0$, que obliga a $x - 3 = 0$ (como en la pregunta 11):
$x = 3$ es la única solución --- la comprobación del
\cref{exo:g8:equations:3} la había encontrado, y la factorización demuestra
que no hay ninguna otra.

\textbf{15.} Por la tercera identidad,
\[
2026^2 - 2025^2 = (2026 - 2025)(2026 + 2025)
= 1 \times 4051 = 4\,051 .
\]
En general $(n+1)^2 - n^2 = (n + 1 - n)(n + 1 + n) = 2n + 1$: la
\omterm{ex:g1:subtraction:difference}{diferencia} de dos cuadrados consecutivos es la \omterm{def:g1:addition:def}{suma} de los dos
números --- el mismo $2n + 1$ de la pregunta 8, obtenido ahora a partir de
la tercera identidad en lugar de la primera.
\end{solution}
